\section{The companion repositories}\label{sec:companion}

The work this paper reports was not all done in one repository. Besides
\texttt{crypto}, where the harness, the ledger and most studies live, the
laboratory kept ECDLP work in nine further repositories
(Table~\ref{tab:repos}). This section says what each holds and what it adds;
results that bear on a family are also stated in that family's section, and
every unit of every repository is listed in
Appendix~\ref{sec:companion-catalogue}. Paths into these repositories are
written \texttt{ar:}, \texttt{ca:} and so on with the prefix of
Table~\ref{tab:repos}, and refer to the pinned commit. Three of the
repositories, like \texttt{crypto} itself, are private; a reader without
access can check only the public ones.

\begin{table}[t]
\centering\small
\caption[]{The laboratory's repositories with ECDLP work, at the commits
read for this paper (24 October 2026). Units: the entries catalogued in
Appendix~\ref{sec:companion-catalogue} (study directories, experiments,
evidence and correction records, knowledge records, or top-level entries).}
\label{tab:repos}
\begin{tabularx}{\linewidth}{@{}l l l r >{\raggedright\arraybackslash}X@{}}
\toprule
Prefix & Repository & Commit & Units & Subject \\
\midrule
--- & \texttt{crypto} (private) & \texttt{5cb553810e} & 1{,}444 & harness, ledger, every study of \S\S\ref{sec:koblitz}--\ref{sec:frontier} \\
ar: & \texttt{crypto-autoresearcher} (public) & \texttt{ab0609ebee} & 4{,}534 & multi-agent ECDLP research with its own claim tiers; experiments, evidence, knowledge records \\
ca: & \texttt{cryptanalysis} (private) & \texttt{c5aa7b39bf} & 467 & discrete-log library and attack suite; a separate $n=83$ pipeline; ECC2K-130 GPU rho \\
ecc: & \texttt{ecc-challenges} (public) & \texttt{86f8a57a5e} & 26 & distributed Pollard rho for the 24 Certicom curves \\
ml: & \texttt{ml-cryptanalysis} (private) & \texttt{58494969fd} & 17 & curve-structure learning on toy ECDLP instances \\
--- & \texttt{ecdlp-prime-field-smallness} (private) & \texttt{4f20be9b36} & 7 & whether any ``small'' subset of $E(\FF_p)$ is addition-compatible \\
--- & \texttt{crypto-lean} (public) & \texttt{148f5dc389} & 12 & Lean 4 library and proof gate shared by the repositories \\
--- & \texttt{ecbench}, \texttt{isogeny-labs}, \texttt{isolab}, \texttt{cryptanalysis-rl} & --- & 49 & a copy of the harness, a website, a benchmark runner, and non-ECDLP reinforcement learning \\
\bottomrule
\end{tabularx}
\end{table}

\paragraph{crypto-autoresearcher.} A multi-agent laboratory with separate
coordinator, executor, reviewer, validator and red-team roles, whose
experiments are frozen specifications with claim tiers and authorised
stages, and whose results are evidence records in a ledger. Of its 2{,}779
experiment specifications, 601 have run output; the rest were approved or
drafted and never run, and are catalogued as such. It holds 793 evidence
records, 427 corrections, about 320 knowledge records (findings,
techniques, open problems) and about 8{,}100 literature records; its
June--August history comes partly from an imported Autolab corpus whose
48\,GB byte mirror is local, not on GitHub. Its ECDLP results, all at toy
claim tiers:
\begin{itemize}
\item \emph{Prime fields, June--August.} An imported Autolab round of
  1 June 2026 solved discrete logarithms end to end with an $m=3$ Semaev
  index calculus and measured a cost exponent of $0.889$ against rho's
  $0.5$, IC/rho growing from $21.5$ to $205.7$ over 10--18 bits
  (\art{ar:experiments/EXP-ALPF-032/analysis.md}). By 18--19 July the
  repository had proved the mixed-volume law of the Semaev family
  (\art{ar:research/THM_BKKMV1.md}) and opened ``Does index calculus beat
  Pollard rho for prime-field ECDLP?'' as its first open problem
  (\art{ar:knowledge/open-problems/KN-OPEN-001.md}).
\item \emph{Binary descent.} msolve's solving degree of the descended $m=3$
  Semaev system is 7 at every $n$ from 11 to 23 (a firm lower bound at
  $n\ge17$, where the run walls on matrix size;
  \art{ar:experiments/EXP-ALBIN-008/analysis.md}). A degree-4 reuse closure
  refutes 400 of 400 unsatisfiable $\Sem3$ descents at $n=19$ where plain
  degree-4 Macaulay refutes none
  (\art{ar:experiments/EXP-CERTBIN-060020/}). A crossbred solver's fitted
  total exponent is $0.891$ $[0.869,0.914]$ against a rho control's $0.498$
  over $n=10$--$24$ (\art{ar:experiments/EXP-ICI-001/analysis.md}). On the 60
  published $\Sem4$ instances of Trimoska et al.\ over $\FF_{2^{15}}$ to
  $\FF_{2^{19}}$, WDSat decides the unsatisfiable $n=19$, $\ell=6$ cell in
  2.37\,s against 124.6\,s for CryptoMiniSat with XOR clauses and 56.7\,s for
  CaDiCaL; no Gröbner engine finished, and there is no rho column
  (\art{ar:experiments/EXP-ICPERF-66fd51/}).
\item \emph{Named binary curves.} At $n=131$ the only Frobenius-stable
  subspace dimensions are 0, 1, 130 and 131, so the orbit trick loses
  57--96 bits for $m=2$--5 (\art{ar:experiments/EXP-FROB-ec08b5/}); the
  descended system exceeds the rho budget by 40.8 bits at $n=131$
  (\art{ar:experiments/EXP-ICPERF-783e9e/}). Conditional on a solving degree
  of at most 4---checked only at $n\le21$---Semaev's 2015 estimate is
  $+49.92$ bits over rho at $n=131$, $+6.15$ at 283, $-37.96$ at 409 and
  $-99.19$ at 571 (\art{ar:experiments/EXP-SEMBIN-992e73/RESULTS.md}); the
  study says nothing in the table is a security statement.
\end{itemize}
The repository's own records correct many of its results (an instrument
tautology, a fabricated justification, superseded headline rows), and the
catalogue follows its corrections.

\paragraph{cryptanalysis.} A C11 discrete-log library (BSGS, rho, kangaroo,
GLV rho, $\FF_p^*$ index calculus, GPU and distributed rho) with bindings
and a Rust attack suite; its README says it contains no ECDLP speed-up. It
holds a second $n=83$ Koblitz pipeline (K600, $a=1$) improved through a
Tier~1 ratchet of rank-probe, closure, index-build and rank-finder changes,
whose latest row is $3.92\times$ a tuned rho end to end, cold
(\S\ref{sec:n83}; \art{ca:experiments/koblitz-n83-rankfinder-20261021/README.md}).
Its \art{ca:ecc2k130/} directory is the ECC2K-130 GPU rho client imported
from \texttt{crypto}, now at 20.08 billion iterations per second on one RTX
PRO 6000 for a table-walk variant that cannot join the live campaign's walk.
Most of its 139 $j=0$ prime-curve directories are fixed-generator $[k]G$
table formats, which compute scalar multiples rather than logarithms and
are out of scope.

\paragraph{The others.} \texttt{ecc-challenges} runs a public
distributed Pollard rho on all 24 Certicom curves, with 338 million
imported ECC2K-130 distinguished points and no curve solved. \texttt{ml-cryptanalysis}
(``curve structural discovery'') recovers a planted automorphism law in
rho cost ($0.572$ against $0.577$) and finds no hardness signal on toy
curves. \texttt{ecdlp-prime-field-smallness} asks whether any
coordinate-defined subset of $E(\FF_p)$ interacts with addition better than
a random subset, and finds none on toy curves---negative evidence on the
prime-field question of \S\ref{sec:prime}. \texttt{crypto-lean} holds the
Lean statements the papers rely on, among them the $\Sem3$ root identity of
Proposition~\ref{prop:s3}. \texttt{ecbench} is a pinned copy of the
\texttt{crypto} harness; \texttt{isogeny-labs}, \texttt{isolab} and
\texttt{cryptanalysis-rl} hold no ECDLP results.

\paragraph{Overlaps.} The ECC2K-130 GPU rho, the PDP oracle programme and
the $n=83$ surfaces study exist in both \texttt{crypto} and
\texttt{cryptanalysis}; the narrative cites the \texttt{crypto} copy. The
two $n=83$ pipelines are separate (K1107 in \texttt{crypto}, K600 here),
and \texttt{crypto}'s validation sweep records that their numbers are not
comparable (\art{research/validation_sweep_20261019/D-known-issues/REPORT.md}).
